\chapter{Types across the boundary}
\label{ch:types}

Reflection tells rosetta that a method takes a \code{GEO::vector<double>}. It
does not tell it what a Python caller should hand over. That mapping --- the
\emph{type gate} --- is per backend, and this chapter is about widening it for
the types your library actually uses.

The default gate handles the obvious cases without a word in the manifest:
arithmetic types, \code{bool}, \code{std::string}, enums, \code{std::vector<T>}
of any of those, and any \emph{bound class}. Everything below is for the rest.

\begin{note}
The IR's rule for anything it cannot claim is uniform: mark the type
\code{kind: "unknown"} and let each backend skip the member. That is why a
type rosetta does not understand costs you one skipped method with a stated
reason, not a compile error.
\end{note}

\section{Foreign sequence containers}
\label{sec:sequences}

Many libraries carry their bulk data in their \emph{own} vector type ---
geogram's \code{GEO::vector<T>}, an aligned or pooled vector --- and the
marshalling layers only know \code{std::vector}. List the container template
(qualified, \textbf{one type parameter}) and it crosses like a
\code{std::vector} of its element:

\begin{lstlisting}[style=json]
"sequences": ["GEO::vector"]
\end{lstlisting}

\code{rosetta\_gen} emits

\begin{lstlisting}[style=cpp]
template <typename T>
struct rosetta::is_sequence<GEO::vector<T>> : std::true_type {};
\end{lstlisting}

into the generated \code{bindings.h}. The container must be
default-constructible with \code{value\_type}, \code{size()}, \code{resize(n)}
and \code{begin()}/\code{end()}; elements must be arithmetic, \code{bool},
\code{std::string} or a bound enum.

The opted-in backends --- \lang{python}, \lang{nanobind}, \lang{node},
\lang{wasm}, \lang{lua}, plus \lang{typescript} declarations --- marshal it
\textbf{by copy through a \code{std::vector<element>} boundary} inside an
emitted adapter. Scripts pass and receive plain arrays, lists or tables. Every
other backend keeps skipping the type.

Three consequences worth knowing:

\begin{itemize}
  \item \textbf{Mutable \code{Seq\&} parameters bind, input-only.} The adapter's
    local container is a real lvalue, so geogram's
    \code{assign\_points(vector<double>\&, dim, steal)} works --- \code{steal}
    steals from the adapter's copy, which is fine. In-place mutations are
    discarded, exactly like pybind's \code{std::vector\&} casters.
  \item \textbf{Overload sets need the sequence overload to be first.} The
    adapter calls the method \emph{by name} with concrete arguments instead of
    spelling the ambiguous \code{\&T::name} member pointer. A backend that keys
    methods by name binds the first-declared overload, so a set whose first
    declaration is the raw-pointer one stays skipped there.
  \item \textbf{Virtual methods naming the container cannot be overridden
    script-side} --- the exact spelling cannot round-trip --- but they still
    bind as callable methods.
\end{itemize}

Sequence-typed public \emph{fields} bind as copying properties in
\lang{python}, \lang{nanobind}, \lang{wasm} and \lang{lua}; \lang{node} skips
them.

\subsection{Registering a concrete type}
\label{sec:concrete}

The string form registers a \emph{template}, and the adapter spells the
container it builds as \code{Namespace::Template<element>}. That composition is
wrong the moment the specialisation carries more than its element:
\code{Eigen::VectorXd} is really \code{Eigen::Matrix<double, -1, 1>}, and the
composed \code{Eigen::Matrix<double>} does not compile.

Register the concrete type instead, spelled exactly as the adapter should write
it:

\begin{lstlisting}[style=json]
"sequences": [{ "type": "Eigen::VectorXd" }]
\end{lstlisting}

which emits a full specialisation --- plus the spelling, as
\code{rosetta::is\_sequence} and \code{rosetta::sequence\_cpp\_name} --- rather
than the partial one.

\code{\{ "template": "GEO::vector" \}} is the long form of a plain string entry.
An object with \emph{both} keys, or neither, is an error --- the two cannot be
told apart by looking at the text, and guessing wrong emits code that does not
compile.

This is also the only way to register a container that is not a template at all.

\section{Foreign matrices}
\label{sec:matrices}

\field{sequences} one dimension up. A 2-D type --- \code{Eigen::MatrixXd}, a
library's own dense grid --- is a sequence in no useful sense, so it needs its
own registration:

\begin{lstlisting}[style=json]
"matrices": [{ "type": "Eigen::MatrixXd" }]
\end{lstlisting}

Entries take the \textbf{same two forms} as \field{sequences}: a plain string
for a one-type-parameter template (\code{"mylib::Grid"}, spelled
\code{mylib::Grid<double>} by the adapter), or \code{\{ "type": "\ldots" \}} for
a concrete type.

The type must be default-constructible with \code{value\_type}, \code{rows()},
\code{cols()}, \code{resize(r, c)} and \code{operator()(i, j)}; the element must
be \textbf{arithmetic} --- a grid of strings has no natural script shape.

The opted-in backends marshal it \textbf{by copy through a
\code{std::vector<std::vector<element>>} boundary} --- an array of row arrays.
\lang{typescript} says \code{number[][]}.

\begin{itemize}
  \item \textbf{Row-major, whatever the matrix stores.} \code{operator()(i, j)}
    is the only access the registration promises, so the boundary is built row
    by row regardless of the type's own storage order.
  \item \textbf{A ragged incoming array is squared off} to the first row's
    length: rows are the outer size, columns the first row's; short rows keep
    whatever \code{resize()} left and extra columns are dropped.
  \item \textbf{A mutable \code{Mat\&} parameter binds input-only}, like a
    sequence one.
\end{itemize}

\section{Foreign-library interop}
\label{sec:interop}

A method taking or returning \code{Eigen::VectorXd} is, to the reflection walk,
a method taking an ordinary class. Every backend used to bind it happily --- and
the call then \textbf{threw at run time}, because no caster for that class was
ever registered.

\field{interop} names foreign libraries whose types the target's binding
framework can already marshal on its own:

\begin{lstlisting}[style=json]
"interop": ["eigen"]
\end{lstlisting}

Recognition is by \textbf{enclosing namespace}, not by a list of type names, so
one line covers everything the library owns: \code{VectorXd}, \code{MatrixXd},
\code{Vector3d}, \code{Map}, \code{Ref}, \code{Block}, \code{Array}, the sparse
types --- including spellings a hand-maintained list would have missed.

\begin{center}
\small
\begin{tabular}{@{}lL{9.4cm}@{}}
\toprule
\textbf{Backend} & \textbf{Behaviour} \\
\midrule
\lang{python} & Emits \code{\#include <pybind11/eigen.h>}; the types bind
  \textbf{natively as numpy arrays}, both directions, matrices included. No
  adapter, no copy where the layout allows a view. \\
\addlinespace
\lang{nanobind} & Same, through \code{<nanobind/eigen/dense.h>}. (Sparse is not
  emitted --- \code{<nanobind/eigen/sparse.h>} costs compile time a rare
  signature does not justify.) \\
\addlinespace
every other backend & \textbf{Skips} the member. Deliberate, and an improvement:
  a skipped method is honest, a bound one that always throws is not. \\
\bottomrule
\end{tabular}
\end{center}

Two practical notes: Eigen must be on the include path of the generated binding
(it already is if \field{user\_include} lists it, which it must anyway for the
bound headers to compile), and the caster header is only emitted when a bound
signature actually names such a type --- so declaring \field{interop} on a
library that never exposes one costs nothing.

\subsection{Reaching the caster-less backends}

\lang{node}, \lang{wasm} and \lang{lua} have no caster to lean on, so on their
own they skip every Eigen-typed member. Give them the flat array by
\textbf{also} registering the concrete type as a sequence:

\begin{lstlisting}[style=json]
"interop":   ["eigen"],
"sequences": [{ "type": "Eigen::VectorXd" }]
\end{lstlisting}

The IR then carries both marks, and each backend picks:

\begin{center}
\small
\begin{tabular}{@{}lL{9.4cm}@{}}
\toprule
\lang{python} / \lang{nanobind} & The \textbf{caster wins} --- still numpy,
  still no copy. A dual-marked type costs them nothing. \\
\lang{node} / \lang{wasm} / \lang{lua} & Bind through the \textbf{sequence
  adapter}: in, out and back through a \code{std::vector<double>} boundary.
  Scripts see a plain array. \\
\lang{typescript} & Declares \code{number[]}, matching what those runtimes hand
  out. \\
\bottomrule
\end{tabular}
\end{center}

\textbf{1-D only.} \code{MatrixXd} has no useful reading as a sequence ---
register it under \field{matrices} instead, which splits the backends exactly
the same way:

\begin{lstlisting}[style=json]
"interop":   ["eigen"],
"sequences": [{ "type": "Eigen::VectorXd" }],
"matrices":  [{ "type": "Eigen::MatrixXd" }]
\end{lstlisting}

covers both shapes, in both directions, on every expanded backend.

\section{Two standard types that need no declaration}

Nothing to write in the manifest for these. They are here because a factory API
tends to name both at once, and until recently either one made a signature
quietly unbindable on the caster-less backends:

\begin{lstlisting}[style=cpp]
std::shared_ptr<Mesh> compile_file(const std::filesystem::path &input);
\end{lstlisting}

\subsection*{\code{std::filesystem::path} marshals as a string}

It is described as \code{kind: "string"} with an \code{is\_path} flag.

\begin{center}
\small
\begin{tabular}{@{}lL{9.4cm}@{}}
\toprule
\lang{python} / \lang{nanobind} & \textbf{Native}, through
  \code{<pybind11/stl/filesystem.h>} / \code{<nanobind/stl/filesystem.h>} ---
  callers pass a \code{str} \emph{or} any \code{os.PathLike}, and a returned
  path comes back as \code{pathlib.Path}. \\
\lang{node} / \lang{wasm} / \lang{lua} & Through the same copy adapter the
  foreign containers use: the boundary declares \code{std::string}, the emitted
  code builds a path from it and calls \code{.string()} on the way back. \\
\lang{typescript} & Declares \code{string}. \\
\bottomrule
\end{tabular}
\end{center}

\subsection*{\code{std::shared\_ptr<T>} crosses in the return direction}

Which is where a factory needs it. The pointee \code{T} must itself be a bound
class.

\begin{center}
\small
\begin{tabular}{@{}lL{9.4cm}@{}}
\toprule
\lang{python} & \code{py::class\_<T, std::shared\_ptr<T>>} --- the holder is
  declared automatically for every pointee the module hands out. \\
\lang{nanobind} & Native, through \code{<nanobind/stl/shared\_ptr.h>}. \\
\lang{wasm} & The class registration gains
  \code{.smart\_ptr<std::shared\_ptr<T>>("T\_sp")}, and only for classes that
  actually travel that way. \\
\lang{lua} & Native --- sol2's own \code{unique\_usertype\_traits} handles it. \\
\lang{node} & The JS object \textbf{adopts} the pointer: a third ownership mode
  in the wrapper, next to ``owns'' and ``aliases a member of a pinned parent''.
  The C++ object lives as long as the JS handle does. \\
\lang{typescript} & Declares the \emph{pointee} (\code{Doc}), not
  \code{shared\_ptr}. \\
\bottomrule
\end{tabular}
\end{center}

Two limits:

\begin{itemize}
  \item \textbf{Returns only.} A \code{shared\_ptr<T>} \emph{parameter} stays
    skipped: converting an incoming script object would mean manufacturing a
    control block over memory the binding does not own --- a double-free waiting
    to happen. Take \code{const T\&}.
  \item \textbf{\lang{node}: untrampolined pointees only.} A class with virtual
    methods is wrapped as its trampoline subclass, and the adopting wrapper
    stores a \code{T*}; reading one as the other is undefined, so such a return
    stays skipped there. The other backends have no such restriction.
\end{itemize}

Both work with a non-copyable \code{T} --- which is the point, since a class
handed out by a factory usually is one.

\section{Deciding what you need}

\begin{center}
\small
\begin{tabular}{@{}L{6.4cm}L{7.0cm}@{}}
\toprule
\textbf{Your signature names\ldots} & \textbf{Write} \\
\midrule
your own \code{Vec<T>} template & \code{"sequences": ["my::Vec"]} \\
\code{Eigen::VectorXd}, Python only & \code{"interop": ["eigen"]} \\
\code{Eigen::VectorXd}, everywhere &
  \code{"interop"} + \code{"sequences": [\{"type": "Eigen::VectorXd"\}]} \\
\code{Eigen::MatrixXd}, everywhere &
  \code{"interop"} + \code{"matrices": [\{"type": "Eigen::MatrixXd"\}]} \\
a 2-D grid of your own & \code{"matrices": ["my::Grid"]} \\
\code{std::filesystem::path} & nothing \\
\code{std::shared\_ptr<Bound>} returned & nothing \\
a \code{std::function} parameter & nothing will help --- the member is skipped,
  with a reason. Add an \field{extensions} wrapper if you need it. \\
\bottomrule
\end{tabular}
\end{center}

When in doubt, generate and read \code{coverage.json}: it names every type that
failed a gate, and which one.
